\documentclass[12pt]{article}
\usepackage[utf8]{inputenc}
\usepackage[T1]{fontenc}
\usepackage{hyperref}
\title{Pataxó Hã-Ha-Hãe Language Resources for K-12 Curriculum Development}
\author{Compiled by Claude Sonnet 5.5 (claude-sonnet-5-5) for Andr\'as Kornai}
\date{}
\begin{document}
\maketitle

\section{Introduction}
Pataxó Hã-Ha-Hãe (ISO 639-3 \texttt{pth}; Glottocode pata1261) is a dormant language of Bahia, Brazil, in the Maxakalian branch of Macro-Jê. It is distinct from the Pataxó community's revived language \emph{Patxôhã}, which is a reconstructed and reinvented variety (begun in 1998) used in Pataxó schools; most of the materials found concern Patxôhã generally and are noted as such. Only one source found addresses the Pataxó Hãhãhãe revival specifically. Caution is needed not to conflate the two.

\section{Primary Educational Resources}
\begin{itemize}
\item Proposal for teaching the language at the indigenous school of the Aldeia Baheta (Pataxó Hãhãhãe revival), Cadernos da Abralin, based on archives, elder interviews and fieldwork (Portuguese). \url{https://cadernos.abralin.org/index.php/cadernos/article/download/253/246}
\item Cleideane Poncada Santana (UFMG, 2016), undergraduate thesis proposing teaching material for Patxôhã at the Pataxó school of Barra Velha (general Patxôhã, not Hãhãhãe-specific). \url{https://www.biblio.fae.ufmg.br/monografias/2016/cleideane%20poncada%20santana.pdf}
\item ``Voos na sabedoria: o ensino do Patxôhã na Escola Estadual Ind\'igena Kijetxawê Zabel\^e'', \emph{Policromias}. \url{https://revistas.ufrj.br/index.php/policromias/article/view/46721}
\end{itemize}

\section{Linguistic and Reference Materials}
\begin{itemize}
\item \emph{Revista Linguística} (UFRJ) article on the Patxôhã revitalization, tracing its 1998 beginnings. \url{https://revistas.ufrj.br/index.php/rl/article/view/10433}
\item Wikipedia, ``Pataxó language''. \url{https://en.wikipedia.org/wiki/Patax%C3%B3_language}
\item UFSB research project on Patxôhã revitalization in the Kaí-Pequi Territory, Prado, Bahia. \url{https://ufsb.edu.br/prosis/images/editais/BAP_2020/Projeto_51.pdf}
\end{itemize}

\section{Religious and Community Resources}
No scripture or dedicated community organization pages were found in search results.

\section{Digital and Audio Resources}
\begin{itemize}
\item ``Kuã Mé'a Aregá: Aprender é brincar'', a digital game for teaching Patxôhã (SBGames). \url{https://sol.sbc.org.br/index.php/sbgames/article/view/32341}
\end{itemize}

\section{Curriculum Development Context}
There are no first-language speakers; any curriculum is revival work. Existing Pataxó school practice with Patxôhã offers models, but the Hãhãhãe community's own revival proposal (Aldeia Baheta) should lead for this ISO code. I could verify only abstracts and listings, not full texts.

\end{document}